% Generated draft: requires mathematical and source review.
\documentclass[11pt]{article}
\usepackage[margin=25mm]{geometry}
\usepackage{fontspec}
\setmainfont{Latin Modern Roman}
\usepackage{amsmath,amssymb,mathtools,unicode-math}
\setmathfont{Latin Modern Math}
\usepackage{enumitem,needspace,xurl,hyperref}
\setlength{\parindent}{0pt}
\setlength{\parskip}{5pt}
\setlength{\emergencystretch}{3em}
\newcommand{\N}{\mathbb{N}}
\newcommand{\Z}{\mathbb{Z}}
\newcommand{\Q}{\mathbb{Q}}
\newcommand{\R}{\mathbb{R}}
\newcommand{\C}{\mathbb{C}}
\newcommand{\F}{\mathbb{F}}
\newcommand{\A}{\mathbb{A}}
\newcommand{\PP}{\mathbb{P}}
\newcommand{\E}{\mathbb{E}}
\newcommand{\eps}{\varepsilon}
\newcommand{\dd}{\,\mathrm{d}}
\begin{document}
\section*{UnsolvedMath ID: 30005321}
\textbf{Source problem code:} catalogue entry
\subsection{Definitions and mathematical statement}
The supplied passage refers to a polynomial \(P_n(q)\) defined in its equation~(1), but that equation is absent from the supplied record. Consequently, the family \(P_n\), its coefficient domain, and its degree are not specified here. The passage also invokes an earlier convention concerning zero coefficients without reproducing it. These omissions prevent a complete, unconditional formulation.

Subject to recovering that definition, let
\[
P_n(q)^r=\sum_{k=0}^{d_{n,r}}a_{n,r,k}q^k,
\qquad d_{n,r}=\deg\bigl(P_n(q)^r\bigr),
\]
for each positive integer \(n\) and \(r\in\{2,3\}\). The extracted question asks whether the coefficients, in increasing order of exponent, have the repeating sign pattern \(+,-,-\). If the omitted convention permits zero coefficients in every sign position, and the omitted definition gives real coefficients, the proposed assertion is precisely
\[
\forall n\in\Z_{>0}\;\forall r\in\{2,3\}\;\forall k\in\Z,
\quad 0\le k\le d_{n,r}\Longrightarrow
\begin{cases}
a_{n,r,k}\ge0,&k\equiv0\pmod3,\\
a_{n,r,k}\le0,&k\equiv1\text{ or }2\pmod3.
\end{cases}
\]
This weak-inequality formulation is conditional on the missing zero convention; it is not a verified replacement for that convention. The case \(r=2\) is attributed to Peter Borwein as the Second Borwein Conjecture, and the case \(r=3\) to Chen Wang as the cubic conjecture. The passage separately describes the Third Borwein Conjecture as involving modulus~\(5\); that distinct conjecture is not the cubic assertion formulated here.
\subsection{Short English statement}
For every positive integer \(n\), the question concerns whether both the square and cube of the source's Borwein polynomial have nonnegative coefficients at exponents divisible by \(3\) and nonpositive coefficients at the other exponents, subject to the source's missing convention about zeros. The defining polynomial formula is also missing, so this remains a conditional statement-only draft. Supplied dated curation labels the second conjecture proved and the cubic conjecture partially proved; these labels have not been independently verified.
\subsection{Sources}
\begin{itemize}
\item Local registry: \texttt{data/unsolvedmath/problems.json}, permanent UnsolvedMath ID \(30005321\), title \emph{Borwein Sign Conjectures}, exact source code \texttt{catalogue entry}. The supplied record carries the status \texttt{partially\_solved} and the editorial label \texttt{corrected\_verified}; neither constitutes independent mathematical verification of this draft.
\item Supplied report citation: \emph{Enumerative Combinatorics (2023)}, Oberwolfach Reports, DOI \texttt{10.4171/owr/2022/57}, pp.~3215--3305. \url{https://doi.org/10.4171/owr/2022/57}. The supplied workshop/proposal year is \(2022\). The passage attributes the second conjecture to Peter Borwein and the cubic conjecture to Chen Wang.
\item Supplied literature source labeled Wang--Krattenthaler: \url{https://arxiv.org/abs/2201.12415}. The record supplies no full bibliographic entry for this link.
\item Supplied literature source labeled original OWR report: \url{https://doi.org/10.4171/owr/2022/57}.
\item Literature assessment dated exactly \texttt{2026-08-22}: the supplied curation states that the first and second Borwein sign conjectures are proved, the cubic statement is only partially proved, and further modulus claims remain conjectural. This is supplied dated curation, not an independently verified literature assessment. The background expressly requests independent verification of time-sensitive status and cited claims.
\item The excerpt refers to bibliography entries [1] and [6], whose full entries are absent. It associates [1] with the Borwein conjectures and an Andrews formula, and [6] with Chen Wang's proof of the first conjecture using saddle point approximation. These references cannot be reconstructed reliably from the supplied record alone.
\end{itemize}
\end{document}
